\vspace{-12pt}
\section{Conclusions}
\vspace{-8pt}

We reproduced the SDR for systematic reviews method by Lee and Sun~\cite{lee2018seed} on all the available CLEF TAR datasets~\cite{kanoulas2017clef,kanoulas2018clef,kanoulas2019clef}.
%
Across all three of these datasets, we found that the 2017 and 2018 datasets share a similar trend in results than to the 2019 dataset. We believe that this is due to topical differences between the datasets and that proper tuning of SDR would result in results that better align with those seen in 2017 and 2018.
%
We also performed several pre-processing steps that revealed that the BOW representation of relevant studies could also share a relatively high commonality of terms compared to the BOC representation. Furthermore, we found that the BOC representation for SDR is generally beneficial and that term weighting generally improves the effectiveness of SDR.
%
We also found that multi-SDR was able to outperform single-SDR consistently. Our results also used an oracle to select the most effective seed studies to compare multi-SDR to single-SDR. This means that the actual gap in effectiveness between single-SDR and multi-SDR may be considerably larger.
%
Finally, in terms of the impact of seed studies on ranking stability, we found that although multi-SDR was able to achieve higher performance than single-SDR, multi-SDR generally had a higher variance in effectiveness.
%

For future work, we believe that deep learning approaches such as BERT and other transformer-based architectures will provide richer document representations that may better discriminate relevant from non-relevant studies. Finally, we believe that the technique used to sample seed studies in the original paper and this reproduction paper may overestimate the actual effectiveness. This is because a seed study is not necessarily a relevant study, and that seed studies may be discarded after the query has been formulated. For this, we suggest that a new collection is required that includes the seed studies that were originally used to formulate the Boolean query, in addition to the studies included in the analysis portion of the systematic review.

%Automating systematic reviews through screening prioritisation reduces the time it takes to create systematic reviews, resulting in improved health care delivery. 
Further investigation into SDR will continue to accelerate systematic review creation, thus increasing and improving evidence-based medicine as a whole.